\section{Evaluation}

To evaluate CrashSimulator we did the following.
We implemented a prototype in 5060 LOC of Python 2 and 2585 LOC of C\footnote{
Line count numbers are according to SLOCCOUNT}.  For a test
environment, we chose a virtual install of Ubuntu Linux 15.04 with a 32 bit kernel on i686 hardware.
We ran the tests on a core i7 system with 8GB of RAM. Our
implementation used {\tt ptrace} to interpose on the running program and
interject anomalous behavior into its execution.
We did not handle threaded programs and disabled ASLR and the kernel's vDSO to make replay
easier.  As far as we are aware, none of these choices influenced the bugs
that CrashSimulator can find.  The only major difference we are aware of 
is its performance (which will be discussed later).
Given appropriate anomaly information, our implementation of CrashSimulator can
test for anomalies that exist in other platforms and environments.  

Using this prototype, we answer the following questions about CrashSimulator:

\begin{enumerate}
   \item{Do environmental bugs exist in widely used applications?}
   \item{Can CrashSimulator identify complex anomalies?}
   \item{What types of bugs could not be found using CrashSimulator?}
   \item{What sorts of errors does CrashSimulator make?}
   \item{Is CrashSimulator able to execute tests in a performant manner?}
\end{enumerate}



%15.04 with several modifications put in place to allow CrashSimulator to operate correctly. First, because
%CrashSimulator relies on an application's memory layout being the same across repeated executions, address space
%layout randomization must be disabled.  Second, the kernel's virtual dynamic shared object feature must be
%disabled.  This feature allows certain kernel data items (timestamps, kernel version information, etc.) to be stored
%in user space.  When the system calls that retrieve this information are made by an application they are instead
%intercepted and handled in user space -- a behavior that would preclude CrashSimulator from being able to interact
%with these calls.  Finally, a 32-bit Linux kernel and supporting environment were chosen.  While CrashSimulator
%could have been implemented on either a 32-bit kernel or a 64-bit kernel such a goal would have resulted in a great
%deal of duplicated effort around dealing with the low level differences between these two kernel versions.
%\emph{I've got a lot more I could talk about here but I'm not sure what info normally goes into this section}


%    \subsection{Limitations - What are CrashSimulator's Technical Limitations}

%The current implementation of CrashSimulator has the following limitations that could be addressed by future work.
%
%\paragraph{Coupling to Architecture}

%As some faults injected by CrashSimulator require low level access to the test system's hardware or operating
%system data structures there exists some degree of coupling between CrashSimulator and these components. One
%area of expansion for CrashSimulator is support for more processor architectures and more operating systems.
%CrashSimulator's test launcher as been designed in such a way that these improvements should be trivial to plug
%in once they have been implemented.
%
%\paragraph{Multi-threaded or Multi-process Applications}
%
%The current implementation cannot correctly replay applications that rely on multi-threading or
%multi-processing.  Is due to two factors.  First, the replay tool must be able to force some pre-determined
%order onto the threads or processes in order to prevent situations where the execution portion of a replay ends
%up requesting system calls in a different order than was recorded in the system call trace the system is
%attempting to replay.  Second, the tool needs to be able to correctly monitor multiple processes. {\tt Ptrace}
%may have appropriate facilities for this task but they were not explored for this version of the tool.
%
%\paragraph{Multi-Platform Tracing Tool}
%
%While there are system call tracing tools available for most of the major operating systems available today
%there is not one tool that works across all of them.  This means that if CrashSimulator is going to be able to
%to work with traces from a given operating system, it will require specific parsing logic for the output of the
%tool used to record the traces.  Additionally, CrashSimulator requires that the tracing tools record all data
%passed into a system call and all data returned from the system call either as a return value or as some block
%of memory written into the calling process's memory.  It is possible that some system call tracing tools do not
%support this level of detail.


\subsection{Do environmental bugs  exist in widely used applications?}

% Describe bug. Why is this attack badddk
% How crash simulator tests
% Show results
% Discuss  What should people think about it

Our first question focuses on whether there are any bugs of the sort that
CrashSimulator can find in popular applications.  After all, if such bugs
do not exist in practice, then there is no need for CrashSimulator.
To test whether such bugs exist in real programs, we 
chose two simple anomalies that could be tested in a large
array of programs.  We chose one file system anomaly and one network anomaly 
that used common API calls and thus had the potential to appear in
many programs.  We ran CrashSimulator on applications from GNU Coreutils and
on popular projects as is indicated by Debian's Popularity Contest\cappos{add
cite}.

\subsubsection{Unexpected File Types} Providing unexpected types of files to
applications.

{\bf Method.}
The first anomaly configured CrashSimulator to identify occurs when an
application has to retrieve and process data from a file.  Linux supports
several ``special'' file types apart from the standard ``regular'' file type.
These include directories, symbolic links, character devices, block devices,
sockets, and First-In-First-Out (FIFO) pipes.  While these special files are can
be operated on programatically, the details around these operations differ
widely from regular files.  For example, a character device that acts as an
interface over a random number generator may be essentially infinite in size.
In this case an application that attempts to read the full contents of a file
before processing them will fail in some way if it attempts to do so on such a
file.  In order to ensure a correct execution applications should examine the
files before they are used in order to ensure they are of an appropriate type.

Identifying these bugs involve changing an application's execution trace to
emulate it being provided a file type it does not expect.  For example, the {\tt
  sed} application, which modifies the contents of a text file according to a
provided command string, could be provided something else such as a symbolic
link, a directory file, or a character device instead.  CrashSimulator
accomplishes this by identifying the calls to {\tt stat()}, {\tt fstat()}, or
{\tt lstat()}  that an application makes to examine the file in question prior
to operating on it and modifying the contents of the structure returned to
indicate that the regular file {\tt sed} is expecting is actually one of the
special filetypes.  If the application responds to this injected information
then there is the possibility that the special file is being handled correctly.
If the application does not alter its behavior then the anomalous condition is
not being handled correctly.

{\bf Findings.}
The results of this work are shown in Table~\ref{table:unexpectedtypes}.  The
condition tested describes for which file type information was modified.
Multiple executions were used to examine the applications response to each of
the listed file types. The ``I...'' column headings denote which file
type was being tested for a given execution.  A result of ``Initial Value''
indicates that this type is the file type the application was provided when the
execution trace being replayed was recorded -- that is, a file of the type the
application was expecting.  A result of ``Recognizes'' indicates that the
application identified that it was being provided with an unexpected file type
and its execution diverged from the trace being replayed indicating that it was
potentially handling the unexpected file type correctly.  A result of ``Fail''
indicates that the application failed to recognize the presence of the unusual
file type because its execution did not diverge from the trace being replayed.


\begin{table*}[t]
    \scriptsize{}
    \begin{tabular}{l  l  |  l  l  l  l  l  l  l}
    \toprule{}
        Application & Condition Tested           & IFREG        & IFDIR        & IFCHR     & IFBLK    & FIFO      & IFLNK    & IFSOCK\\
        Aspell      & Dictionary File            & Initial Value  & Fail           & Recognizes  & Fail       & Fail        & Fail       & Fail\\
        Aspell      & File being checked         & Initial Value  & Fail           & Recognizes  & Fail       & Fail        & Fail       & Fail\\
        gnu-gpg     & secring.gpg                & Initial Value  & Fail           & Fail        & Fail       & Fail        & Fail       & Fail\\
        vim         & File being opened          & Initial Value  & Recognizes     & Recognizes  & Recognizes & Recognizes* & Recognizes & Fail\\
        nano        & File being opened          & Initial Value  & Recognizes     & Recognizes  & Recognizes & Fail        & Fail       & Fail\\
        sed         & File being edited          & Initial Value  & Fail           & Recognizes  & Fail       & Fail        & Fail       & Fail\\
        df          & /proc                      & Fail           & Initial Value  & Fail        & Fail       & Fail        & Fail       & Fail\\
        wc          & File being checked         & Initial Value  & Recognizes     & Recognizes  & Recognizes & Recognizes  & Recognizes & Recognizes\\
        du          & Directory being checked    & Recognizes     & Initial Value  & Recognizes  & Recognizes & Recognizes  & Recognizes & Recognizes\\
        install     & File being installed       & Initial Value  & Recognizes     & Fails       & Fails      & Fails       & Recognizes & Fails\\
        fmt         & File being formatted       & Initial Value  & Fails          & Recognizes  & Fails      & Fails       & Fails      & Fails\\
        od          & File being dumped          & Initial Value  & Fails          & Recognizes  & Fails      & Fails       & Fails      & Fails\\
        ptx         & File being read            & Initial Value  & Recognizes     & Recognizes  & Recognizes & Recognizes  & Recognizes & Recognizes\\
        comm        & Second file being compared & Initial Value  & Fail           & Recognizes  & Fail       & Fail        & Fail       & Fail\\
        pr          & File being read            & Initial Value  & Fail           & Fail        & Fail       & Fail        & Fail       & Fail\\
        readlink    & XXXXXXXXXXXXXXXX           & XXXXXXXX       & XXXXXXX        & XXXXXXXXXX  & XXXXX      & XXXXXXXXX   & XXXXXXXXX  & XXXXXX\\
        mkdir       & XXXXXXXXXXXXXXXX           & XXXXXXXX       & XXXXXXX        & XXXXXXXXXX  & XXXXX      & XXXXXXXXX   & XXXXXXXXX  & XXXXXX\\
        link        & XXXXXXXXXXXXXXXX           & XXXXXXXX       & XXXXXXX        & XXXXXXXXXX  & XXXXX      & XXXXXXXXX   & XXXXXXXXX  & XXXXXX\\
        mknod       & XXXXXXXXXXXXXXXX           & XXXXXXXX       & XXXXXXX        & XXXXXXXXXX  & XXXXX      & XXXXXXXXX   & XXXXXXXXX  & XXXXXX\\
        truncate    & XXXXXXXXXXXXXXXX           & XXXXXXXX       & XXXXXXX        & XXXXXXXXXX  & XXXXX      & XXXXXXXXX   & XXXXXXXXX  & XXXXXX\\
        unlink      & XXXXXXXXXXXXXXXX           & XXXXXXXX       & XXXXXXX        & XXXXXXXXXX  & XXXXX      & XXXXXXXXX   & XXXXXXXXX  & XXXXXX\\
    \bottomrule{}
    \end{tabular}
    \caption{Applications tested for their handling of unexpected file types.}
    \label{table:unexpectedtypes}
\end{table*}

The frequency of failed executions in our results indicate that many
applications make the assumption that they will only be used to process
regular files.  When this assumption does not hold the results off the
application's execution are hard to predict.  In many cases a denial of
service condition occurs in the form of the application ``hanging'' as it
attempts to incorrectly process the file.  This may happen harmlessly (such
as the case where an application blocks forever waiting for a {\tt read()}
call to retrieve non-existant data from an empty FIFO, or harmfully, as
happens in the case where an attempts to read in and process an
``infinitely large'' file eventually filling available memory or disk
space all the while consuming as much processing power as the operating
system can give it.  Either of these outcomes would be particulary
detrimental if the application in question is run as part of an unattended
scripted process or workflow.

\preston{I need to collect these results}
An interesting behavior we recognized when analyzing these results can be
seen in the bottom 6 applications listed in the above table.  Each of these
applications correctly recognized when any filetype other than its expected
filetype.  This is because these applications are simply ``dump'' wrappers
around they system calls they share a name with.  Essentially, they just
hand the specified file to their corresponding system call and, in the case
of an error, use {\\tt perror()} to print the correct localized error
message.  In short, these applications pass their error handling
responsibilities to the kernel.

\subsubsection{Slowloris attack} Delaying network applications for extended
periods.

% Motivation - mentioned and addressed in other areas.  Effort in place to make package managers resiliant to this
% so that you can't be tricked into using old software
% All of these tools failed but we would expect that modern software like apt would be resiliant to it. We were unable
% to test several prominiant network programs because our tool cannot replay multi-threaded or multi-process programs


\cappos{I need to trim this *a lot*}
% Sentence at the end that says "now that we've seen that these issues are prevalent across major applications...

{\bf Method.}
CrashSimulator is not limited to detecting filesystem-based bugs.  Another area
in which CrashSimulator has shown promising results is identifying bugs around
network communications.  The second anomaly we examined using CrashSimulator
involves an applications behavior when it attempts to communicate over a network
with extremely long (on the order of minutes) response times.  At a low level,
applications tend to follow a similar pattern when retrieving data from a
network socket.  Specifically, they use some sort of a looping construct repeat
two actions: they wait for data to be available from the network with using one
of a few system calls meant for this purpose and then they read the data
into a userspace buffer using a system call.  A key aspect of this approach is
detecting when the user should be notified of an issue because too much time has
passed during the ``waiting'' step.  This takes the form of a timeout value
which indicates after how much time the ``waiting'' operations should result in
an error code that the application can handle.

At a system call level these two steps are usually implemented as follows.  The
waiting process is mediated by a call to {\tt poll()} or {\tt select()}.  These
system calls return when one or more of the file descriptors they are monitoring
are ready to be read from (or written to).  These system calls support a timeout
value passed as a parameter.  After this timeout value has passed they return an
error giving the application an opportunity to take action.  Alternatively, an
application might simply use {\tt read()} or {\tt receive()} in a loop to
retrieve all data from a socket.  In this case, a well behaved application
should specify a timeout value for the socket itself to get the same
functionality as described in the former case.

This approach falls apart in situations where developers specify timeout values
that are much longer than is reasonable.  This is an issue that CrashSimulator
can expose.  As an example, consider {\tt wget}.  Its primary data reception
loop relies on a {\tt select()} call to identify when data is available to be
read.  This {\tt select()} call is given a timeout value of 899 seconds -- a
value one second short of Linux's default maximum TCP connection timeout value.
It can be argued that if something close to this timeout value occurs during
even one iteration of {\tt wget's} reception loop the user has experienced an
adverse execution. Taking this to the extreme, situation where every iteration
of the loop takes just under the maximum timeout value is clearly a bad
execution.  A well behaved application should specify a reasonable timeout value
after which it takes some action to recover from the poor network conditions or,
at the very least, alerts the user that an unusual situation has occurred.

As in the previous unusual file types example, CrashSimulator can make claims
about whether an application is vulnerable to this sort of situation by
modifying injecting anomalous data into the results of these system calls --
namely large timeout values -- and observing whether or not the application
takes any action.  Addionally, by properly manipulating the results of all
time-returning calls, CrashSimulator can simulate an execution where close to
the maximum timeout value occurred in each iteration nearly instantly.

Additionally, CrashSimulator can passively analyze a replay execution of an
application performing network communication in order to determine whether it is
vulnerable to this sort of attack.  In this approach CrashSimulator determines
whether or not the application makes any effort to configure its network
communications with a timeout value. This is done by examining the presence or
absence of {\tt setsockopt()}, {\tt poll()} and {\tt select()} calls as well as
the timeout values that may or may not have been passed to them. Applications
that do not do so are vulnerable to this attack by default as the operating
system defined protocol timeout value is used as a fallback.

We use CrashSimulator to analyze 10 network applications and libraries in
widespread useage (based on Debian's popularity contest ratings).  The results
of this analysis can be found in Table~\ref{table:slowloris}.  The analysis results in
this table indicate whether or not the applications where vulnerable and, in the
case of vulnerablilty, a brief description of the cause.

\begin{table*}[t]
  % Try apt-get, they have reported this as a bug
  \scriptsize{}
  \begin{tabular}{l | l}
    \toprule{}
    Application              & Analysis Result\\
    wget                     & Overly long timeout supplied to {\tt select()}, vulnerable\\
    ftp                      & No {\tt poll()} or {\tt select()}, no timeout set, vulnerable\\
    telnet                   & {\tt select()} specifies no timeout, vulnerable\\
    urllib http              & No {\tt poll()} or {\tt select()}, no timeout set, vulnerable\\
    urllib ftp               & No {\tt poll()} or {\tt select()}, no timeout set, vulnerable\\
    ftplib                   & No {\tt poll()} or {\tt select()}, no timeout set, vulnerable\\
    httplib                  & No {\tt poll()} or {\tt select()}, no timeout set, vulnerable\\
    requests                 & No {\tt poll()} or {\tt select()}, no timeout set, vulnerable\\
    urllib3                  & No {\tt poll()} or {\tt select()}, no timeout set, vulnerable\\
    python-websocket-client  & No {\tt poll()} or {\tt select()}, no timeout set, vulnerable\\
    \bottomrule{}
  \end{tabular}
  \caption{Applications tested for their handling of extremely slow response
    times from the host they are communicating with}
  \label{table:slowloris}
\end{table*}


{\bf Findings.}
As can be seen from the above results, all of the applications we examined were
vulnerable to this sort of anomalous condition.  What's more, in the vast
majority of cases it is because the application makes no effort to specify a
timeout value.  This means that the application is effectively falling back on
the operating system's maximum timeout value as defined for the protocol in
use.  This value is \emph{19 minutes} for the above applications as defined by
Linux for TCP sockets.

This situation has effectively exploited in the wild.  The {\tt slowloris} tool
allowed a single attacking host to successfully denial of service a vulnerable
web server with minimal bandwidth usage.  This attack worked by introducing
delays between the time headers were sent to the victim by an attacking HTTP
client.  Because many web servers at the time allowed large timeouts between
messages from the clients they were communicating with, one web client could
effectively tie up the resources of a web server for a long period of time.

As discussed above, the version of wget we examined is vulnerable to this sort
of attack from the other direction.  Consider a malicious web server that
monitors connecting clients with the goal of specifically targeting {\tt wget}.
When it sees a connection from {\tt wget} it begins a reply process where
messages are only sent every 800 seconds (remember that {\tt wget's} default
timeout value is 899 seconds.  If {\tt wget} is not being actively monitored its
transaction would take hours to complete.  This is particularly dangerous in
situations where {\tt wget} is being run as a scheduled job.  It is very likely
that future iterations of the job could be kicked off before the intitial job
completes.  At best there are now multiple iterations of the job conflicting
with one another.  At worst, there could be a slow buildup of unending or long
lived jobs consuming resources and causing deterimental system behavior.
 \cappos{cite SlowLoris}

Because this attack has been seen in the wild, the developers of major projects
have taken steps to ensure their applications are resiliant against it.  The
{\tt apt} project is one such group.  We were unable to examine {\tt apt} and
several other major neetwork applications because of the previlence of multiple
threads and porcesses as a common pattern in handling network communications.
We expect that CrashSimulator would find these major applications sufficiently
protected.


{\bf Concluding Thoughts}
The results of the above analysis show that CrashSimulator is readily able to
identify novel bugs in existing applications with mature code bases.  This is
because CrashSimulator focuses on identifing bugs that result from anomalous
conditions present in an applications environment.  For the unexpected file
types bugs the fact that the anomalous conditions were handled in some cases but
not in others indicate that an understanding that unexpected file types may be
encountered is present in the applications' development communities but also
that this understanding is incomplete.  The presence of these bugs in major
security-centric applications like {\tt gnu-gpg} is of particular concern given
the scrutany that its codebase has seen.  The fact that none of the applications
we looked at in our ``slowloris'' analysis correctly dealt with the situation
indicates that either the devleopment communities around these applications are
unaware of the possiblilties of such situations or that they do not believe
these situations are dangerous.  As we have discussed above this is not the
case.


%Thus we see that even very
%widely used applications are subject to failures due to unexpected file
%types.  The only major exception is those applications like {\tt mkdir},
%{\tt ls}, and {\tt chown}, that are thin wrappers on system calls.  This
%is because the operating system kernel handles error handling for those
%calls instead of the application

\subsection{Can CrashSimulator identify complex anomalies?}

% Re-organize this to the same format as the above two sections
% Talk about doing manual source code analysis to identify interesting edge conditions

Given that now we see that many applications have simple environmental bugs,
we wanted to understand if CrashSimulator can identify more complex types
of misbehavior that crosses many system calls.  To do this, we examined the
way in which applications move a file.  While this can be done via the {\tt
rename()} system call in some cases, moving files is far more complex in
many cases.  For example, the file must be copied if moved across file
systems, which adds additional complexity.  To come up with the correct
system call trace, we examined the system calls made by the {\tt mv}
command and compared the behavior of other tools.  We found that other
tools made many errors when trying to move files.

In Linux, the rename() system call will only move a file if the source and
destination are on the same device.  This means that moving a file from a
directory structure location on one storage device to a directory structure
location on another device must be handled on a case-by-case basis by any
application that relies on this operation.  

{\bf Method.}  We created anomalies for the move operation as follows.
We examined several different move libraries and found that {\tt mv} seemed
to handle cases that other tools failed to consider.  Following this, we 
used the steps of {\tt mv} as a template and
constructed computational models \cappos{What is this?  I'm confused by what
you are saying here.} that can determine whether or not a replayed
execution of an application has performed all the necessary steps to
successfully carry out a cross device move.  This analysis included manual
review of system call traces of {\tt mv} moving a file across devices as well as
analysis of its source code for interesting checks and edge conditions.  Each of
the following are a ``step'' in the process of moving a file across devices that
a well behaved application should carry out.


% Removed based on discussion
%\paragraph{Verify Destination is not Target of Source}
%
%In this condition, CrashSimulator confirms whether or not the application under test performs a check to ensure
%that the source file name is not a symlink pointing to the destination file name.  If this check is not
%performed, loss of data is possible due to removal or overwriting of the destination during the copy process
%resulting in the source symlink pointing to nothing. \emph{!!!! This is likely one that we would have trouble

%    detecting based only on system call behavior}
\paragraph{Ensure Source is not Replaced Between Check and Move - Source Not Replaced}

The first condition we configured CrashSimulator to look for involves ensuring
that the file being copied is not replaced during the move process. Not
performing this step is an example of a classic race condition.  An application
must make an effort to ensure that the file being copied is not replaced between
the time it is initially examined and the time it is opened for copying.  {\tt
  mv} handles this situation by storing the inode number returned by an initial
{\tt stat()} call and compares it to the inode number returned by an {\tt
  fstat()} call made after the file has been opened for copying.  If the file
has been replaced on the filesystem these inode numbers will differ and {\tt mv}
will alert the user to the situation.  This strategy works because Linux does
not truely remove a file until all applications that have a file descriptor to
the file close it.  This means that if the file is intact after it has been
opened, it is guaranteed to remain intact until the application finishes with it
and calls {\tt close()} on its file descriptor. CrashSimulator is able to
identify situations where the above saftey checks do not happen by monitoring a
replay execution of a system call trace of an application moving a file across
devices.  If the requisite system calls are not present then the application has
not performed the file move safely and the execution is rejected.

\paragraph{Preserve Extended File Attributes - Preserve Xattrs}

The second condition we configured CrashSimulator to identify is whether or not
an application properly preserves extended file attributes when moving a file
from on device to another.  Extended file attributes are used by modern
operating systems to store descriptive information about a file that cannot be
held by normal filesystem fields.  For example, an operating system can use
extended file attributes to store whether or not a file was downloaded from the
internet -- information that can be later used in a security context to warn
users if they are about to open a potentially unsafe file.  {\tt mv} handles
preservation of these attributes by making the appropriate system calls to read
all of the extended file attributes off of the source file and then applying
them to the destination file.  Once again, CrashSimulator is able to determine
whether or not an application has correctly preserved these attributes by
monitoring a replay execution of an application moving a file across devices.
Specifically, CrashSimulator watches for the application to make system calls
that read the extended file attributes from the source file (i.e. {\tt
  getxattr()}, {\tt lgetxattr()}, or {\tt fgetxattr()}) followed by system calls
that re-apply the attributes to the destination file (i.e. {\tt setxattr()},
{\tt lsetxattr()}, or {\tt fsetxattr()}).

\paragraph{Preserve Timestamps - Preserve Timestamps}

The third condition we configured CrashSimulator analyze the preservations of
timestamps from the source file. In order to truely copy a file from one device
to another an application needs to ensure that its metadata as well as its
contents are preserved in the destination file.  {\tt mv} ensures that this
occurs by making system calls that retrieve the POSIX timestamps associated with
a file (content modification time, file status change time, and last access
time) and apply them to the destination file.  CrashSimulator is able to monitor
a replay execution in order to ascertain whether or not an application applies
the appropriate timestamps to the destination file as part of moving a file
across devices by monitoring for a system call from the {\tt stat()} family to
retrieve the timestamps and system call from the {\tt utime()} family to apply
the retrieved timestamps to the destination file.

\paragraph{Move Infinite Character Device Across Disks - Careful Copying Character}

The final condtion we configured CrashSimulator to diagnose is whether or not an
application identifies the situation where it has been instructed to move a
special character file like {\tt /dev/urandom} across disks.  As was discussed
in the previous section, many applications fail to examine the nature of a file
before operating on it.  This is one case where {\tt mv} does perform such
checks.  It identifies the situation and correctly moves the special file rather
than engaging in a manual, cross-disk copy process.  If an application fails to
detect this situation, the typical steps in a manual copy process could result
in the application filling the destination device, consuming available memory,
or simply hanging indefinitely as it attempts to read in the full contents of an
effectively infinite-sized file.  CrashSimulator is able to detect this
mis-behavior in two ways.  The first is to replicate the steps in discussed in
the previous section on unexpected file types with the goal of injecting data
that indicates that the source file is a character device.  The second is to
monitor a trace where the application actually attempts to move one of these
special files and identifying whether the application executes a series of {\tt
  open()}s, {\tt read()}s and {\tt writes()} that indicate it is incorrectly
attempting to exhaustively read the file.

Table~\ref{table:crossdevice} lists the result of our analysis of several applications
and libraries that purport to correctly move files within a filesystem.  The
column headings correspond to the short names listed for each of the previous
paragraphs. A result of ``Correct'' indicates that the application or library
function correctly performs the step described.  A result of ``Incorrect''
indicates that the application failed to perform the step correctly or didn't
make an effort to perform it at all.

 \begin{table}[t]
    \scriptsize{}
    \begin{tabular}{l l l l l}
    \toprule{}
        Application     & Source Not Replaced & Preserve Xattrs & Preserve Timestamps & Careful Copying Character\\
        mv              & Correct             & Correct         & Correct             & Correct\\
        mmv             & Correct             & Wrong           & Wrong               & Correct\\
        install         & Correct             & Wrong           & Wrong               & Wrong\\
        perl File::Copy & Correct             & Wrong           & Wrong               & Wrong\\
        shutils         & Wrong               & Wrong           & Correct             & Correct\\
        rust            & Correct             & Wrong           & Wrong               & Wrong\\
        boost::copyfile & Wrong               & Wrong           & Wrong               & Wrong\\
    \bottomrule{}
    \end{tabular}
    \caption{Applications and libraries analyzed to determine whether or not
      they are able to correctly move a file from one device to another}
    \label{table:crossdevice}
\end{table}

{\bf Findings.}
As can be seen from the results in Table~\ref{table:crossdevice} each of the
applications tested fails to perform one or more of the steps required to
successfully complete a cross-device move.  This is an unfortunate situation
because a failure to perform any one of these steps can result in negative
outcomes for the system as a whole.  If the ``Source Not Replaced'' step isn't
performed, it cannot be said for certain that the intended data will make it
into the destination file.  The same can be said for the ``Preserve Xattrs''
step.  This can result in security failures as real world systems rely on the
information contained in extended file attributes to reason about the safety
around a given file.  For example, Apple's Gatekeeper relies on extended file
attributes to prevent applications downloaded untrusted developers from being
executed without user.  Additionally, failure to preserve timestamps has obvious
implications in records-keeping and archival situations as well as any situation
where a chronological record of file maniuplations is important.

The most dramatic failures we saw happened when an application failed to perform
the ``Careful Copying Character'' checks.  Applications that blindly copied from
an infinite character device filled both disks and memory while consuming
maximum computing resources

Our results indicate that CrashSimulator is able to identify whether complex
operations are performed correctly or not.  These analysis rely on
CrashSimulator's ability to compose sets of checkers that each deal with a step
from the operation.  We have real world validation of this in the form of a bug
report found on Python's bug tracker after we completed our analysis.  Python
bug 15100 reports issues around the {\tt shutils} module's functions for copying files
when they must do so across disks.  The report states the module doesn't do
sufficient checking to prevent the race condition described in ``Source Not
Replaced''.CrashSimulator was able to correctly identify both that {\tt shutils} was
not correctly performing the move operation and diagnose specifically that the
``Source Not Replaced'' race condition was present.

\subsection{What types of bugs could not be found using CrashSimulator?}

Due to the level at which CrashSimultaor works (the system call level), it is
not possible to find certain types of bugs.  CrashSimulator is not able to
detect bugs that do not result from or have a visible effect on the system calls
an application makes.  An example of this situation exposed during work on this
evaluation is applications mishandling a situation around file moves where the
where, due to a strange volume mounting arrangement, the source directory is
moved into itself.  In order to detect this situation the application must make
an effort to gather the full path of both the source and destination and compare
them.  While CrashSimulator is able to observe some aspects of the ``gathering''
process (i.e. the {\tt stat()} or {\tt lstat()} calls) it is not able to observe
whether a proper comparison takes place.  As a result, the best a CrashSimulator
checker could do in this situation is observe a lack of gathering-related system
calls and issue a warning that the issue might be present.

%\emph{CrashSimulator could
%  detect this issue during a specially crafted execution which replays a trace
%  recorded where the issue was present.  I'm not sure how useful this would be
%  as presumably the while recording the trace the user would be able to directly
%  observe whether things happened correctly or not}
        
\subsection{What sorts of errors does CrashSimulator make?}
% Talk about both false positives and false negatives
% When and how they occur
% False negative would be checkers that aren't specific enough, or a given test case/checker is not included
% OR the system call trace does not include all of the information you need to know if something has happened
% Keep this short
There are two types of errors that CrashSimulator makes. First, it can 
incorrectly flag a situation as being an error (a false positive).  Second,
it can fail to correctly detect an issue that actually occurs (a false 
negative).  We discuss each situation in turn.

{\bf False Positives.}
The primary source of false positives (that is, CrashSimulator reports a failure
in a situation where the execution should be accepted) for CrashSimulator is an
application using a different sequence of system calls than a checker was
anticipating to implement a given operation.  Fortuantely, CrashSimulator's
approach allows these situations to be easily corrected once identified.

In many cases a given operation can performed in several ways.  This means that
a given operation may be correctly implemented by a set of system call sequences
that differ from each other on a variety of dimensions.  In order to correctly
classify applications that make use of alternative methods for a given operation
the compuational models CrashSimulator relies on must accept the assoc

As an example consider GNOME's {\tt glib} file handling facilties.  When an
application makes use of these facilities to move a file across storage devices
the library itself correctly handles the above behaviors we definied earlier as
making up a "correct" file move operation.  When we used CrashSimulator to
analyze a minimal application using {\tt glib} to perform this operation we
initially received failing results.  We found through manual examination of a
the system call trace used in this analysis that, while {\tt glib} correctly performs
the required behaviors, it does so using alternative system call sequences.  For
example, the typical way to move the actual contents of a file from one file to
another is to {\tt read()} the data from the source file and {\tt write()} the
data back out to the destination file.  Glib instead creates a pipe and uses the
{\tt splice()} system call to copy the contents out of the source file, through
the pipe, and into the destination file.

This approach is viable because there are a finite number of system calls, and a
given operation can be mapped to a managable subset of system calls.  Given the
above example around moving files, consider the mapping from high level
``operation'' to the set of system calls that can implement it in
Table~\ref{table:stepsandcalls}.  Each of the steps in the operation map to a small
number of system calls and, in most cases, only one of the system calls is
necessary.  In situations where two system call sequences can correctly
implement the same operation CrashSimulator simply runs two computational models
in parallel and accepts the execution if either model ends in an accepting
condition.

\begin{table*}[t]
    \scriptsize{}
    \begin{tabular}{l | l }
    \toprule{}
      Operation                                               & Potential System calls\\
      Examine source file                                     & stat64(), lstat64(), fstat64()\\
      Examine destination file                                & stat64(), lstat64(), fstat64()\\
      Open source file                                        & open()\\
      Read contents of source file                            & read(), splice() with a pipe\\
      List source file's extended file attributes             & listxattr(), llistxattr(), flistxattr()\\
      Read contents of source file's extended file attributes & getxattr(), lgetxattr(), fgetxattr()\\
      Open destination file                                   & open(), optionally unlink() the file first\\
      Write contents to destination file                      & write(), splice() with a pipe\\
      Apply extended file attributes to destination file      & setxattr(), lsetxattr(), fsetxattr()\\
      Apply proper timestamps to destination file             & utimens(), futimens()\\
      Apply proper permissions to destination file            & chmod(), open() with a modeline specified\\
      Close the source file                                   & close()\\
      Close the destination file                              & close()\\
    \bottomrule{}
    \end{tabular}
    \caption{Each step of a successful cross-disk file move operation mapped to
      the system call or calls that can implement it}
    \label{table:stepsandcalls}
\end{table*}
    
{\bf False Negatives.}
CrashSimulator can encounter false negatives (that is, CrashSimulator accepts an
execution that it shouldn't) in three situations.  The first is that is hasn't
been provided with the configuration necessary to test for the issue it missed.
This is essentially the same as a ``missing testcase'' sort of situation that is
common throughout testing methodologies.  CrashSimulator cannot identify issues
it doesn't know how to look for.  The second situation is when a computational
model is not constructed correctly and, as a result, misses an issue it should
identify.  The third case, which relates to the first two, involves a situation
where the trace level is inadequate to get enough semantic information to 
make a correct judgment.  For example, system call tracing does not provide
insight into computational actions that may be performed, such as an 
application deciding to cache information in memory when it fails to open a 
file.

\subsection{Is CrashSimulator able to execute tests in a performant manner?}

\cappos{All of the performance numbers should be here.  I recommend having
a graph with the normal run time on the x-axis and the replay speed on the
y-axis.  It will be a scatter plot.  Each point will be one of the
experiments you ran.}

One key attribute of successful testing tools is that they be able to complete their tests in a timely manner.
If a tool takes too long to complete its tests users will be less likely to run it frequently, or at all,
reducing the tools overall usefulness dramatically. To this end, the performance of CrashSimulator was evaluated
in order to determine whether or not it was able to complete its test executions in an acceptable time frame.

Each of the following evaluations examines the completion times for 20 consecutive executions of the specified
application in both native and replay execution modes.  Multiple consecutive executions were used for two
reasons.  First, a single execution in either mode completes so quickly for the applications in question as to
be nearly indistinguishable from the other mode.  Second, multiple executions somewhat contros for varying
conditions like filesystem activity or network latency that can vary between individual executions.


    \begin{figure}[t]
        \center{}
        \includegraphics[scale=.5]{FakeTimeEval}
        \caption{\emph{Fake data that shows the runtime differences with
CrashSimulator vs the native application.}}

    \end{figure}


    \begin{table}[H]
        \scriptsize{}
        \begin{tabular}{l  l  l  l}
            \toprule{}
                Execution Description & Native Eecution & Replay Execution\\
                network\_speedtest & 0.024 & 2.407 \\
                filesystem\_speedtest & 0.038 & 2.604 \\
                mv (cross-disk file move) & 0.016 & 2.995 \\
            \bottomrule{}
        \end{tabular}
        \caption{This sum total runtime of 20 runs of the listed application.}
    \end{table}


\cappos{I think these likely all get removed...  I don't understand the purpose}
\paragraph{Discussion on network\_speedtest}

Sample application opens a TCP connection to an already running listener (netcat), sends a message, and
exits. For simple applications like this, native execution is significantly faster than replay execution...

\paragraph{Discussion on filesystem\_speedtest}

This sample application create a new file with open(), writes a message to it, closes the file descriptor, and
exit()'s. This test is particularly meaningful as it replicates a pattern of system calls that is used for in
all sorts of applications in Linux...

\paragraph{Discussion on mv}

In addition to the newly constructed sample programs, performance values for replay of an execution of {\tt
  mv} moving a file between two separate disks were recorded.  Results indicate that CrashSimulator is able to
replay this operation in a similar time period to the time in which it is able to replay the dramatically
simpler programs.  The fact that {\tt mv} makes an order of magnitude more system calls did not substantially
increase replay time.

{\bf Findings.}
Overall, the performance of our CrashSimulator simulator is about two orders of
magnitude slower than executing the original program natively.  This is 
largely due to performance problems in our (unoptimized) prototype, such as
our utilization of {\tt ptrace} instead of more efficient mechanisms like
{\tt eBPF} or a loadable kernel module.  With a more optimized
implementation, we likely could be even more efficient than executing the
commands since we only replay the command system calls and do not need to
do other related operations such as I/O.
